\section{Asian and scope of the instrument}

This appendix reproduces the research instruments of the Adaptive Execution and
Error-Recovery Engine (AEERE, Component~3): the error taxonomy and error list,
the data set design, the noise-injection protocol, the three safety checks, the
verification targets, and the public API.

\section{Error taxonomy and the error-type lists}

The engine keeps two error lists. Step 1 of the labelling cascade uses a list of
eight hand-written error types made just for Kali offensive-security tool
failures, each encoded by hand-written patterns and documented with its fix,
retry limit and wait rule (Table~4 in the component report). The full test data
set extends this to \textbf{ten} error types, adding the two additional classes
distinguished only by the text model and the AI step. The ten types are
\emph{service\_error}, \emph{file\_not\_found}, \emph{port\_conflict},
\emph{syntax\_error}, \emph{dependency\_error}, \emph{network\_error},
\emph{permission\_error}, \emph{authentication\_error}, \emph{package\_not\_found},
and \emph{command\_not\_found}. Each type carries a default retry limit and a
tool-swap rule; retries stop at the limit for the type, and a swap rule decides
when a different tool is the right fix (for example gobuster may be swapped for
ffuf or dirb for folder scanning, and nmap for masscan or rustscan for network
scanning).

\section{The 8,000-case data set}

The training and test data set contains 8,000 labelled command-failure cases.
Each case records the command, the raw error message, the exit code, the correct
error type, the cause, and the best fix. The case count per error type is shown
below; the set covers more than eighteen tools (including netcat, gobuster,
sqlmap, hashcat, nikto, ssh, nmap, docker, dirsearch, and john) and three
difficulty levels.

\begin{center}
\begin{tabular}{@{}lcc@{}}
  \toprule
  \textbf{Error type} & \textbf{Cases} & \textbf{Share}\\
  \midrule
  service\_error & 839 & 10.5\\
  file\_not\_found & 811 & 10.1\\
  port\_conflict & 812 & 10.2\\
  syntax\_error & 812 & 10.2\\
  dependency\_error & 807 & 10.1\\
  network\_error & 806 & 10.1\\
  permission\_error & 802 & 10.0\\
  authentication\_error & 797 & 10.0\\
  package\_not\_found & 759 & 9.5\\
  command\_not\_found & 755 & 9.4\\
  \midrule
  Total & 8,000 & 100.0\\
  \bottomrule
\end{tabular}
\end{center}

All addresses in the cases are private RFC-1918 lab addresses. The set is split
into training (70\%, 5,600 cases) and test (30\%, 2,400 cases) sets, stratified
by error type, under a fixed random seed so the measurement can be repeated.

\section{Noise-injection protocol}

To test whether the labeller generalises beyond clean, hand-typed errors, each
test message is degraded with controlled noise that is realistic for a real
terminal: log prefixes and timestamps, trailing text, changed letter case,
keyboard-style typos, and truncation. The noise level is raised step by step
from 0\% to 60\%. The text model is only ever trained on the clean training
half, so any degradation in the noised test is an honest measure of
generalisation.

\section{The three safety checks}

Every fixed command must pass all three checks before it is allowed to run:

\begin{enumerate}
  \item \textbf{Format check} -- the AI's reply must parse as the expected JSON
  plan.
  \item \textbf{Dangerous-word check} -- blocks destructive or
  fetch-and-execute patterns (for example \texttt{rm -rf}, disk wiping, or
  downloading and running code).
  \item \textbf{Confidence-and-scope check} -- requires a minimum confidence
  level and rejects any target outside private RFC-1918 address ranges.
\end{enumerate}

Every blocked case is written to a log in an SQLite database, which provides
the evidence base for the dangerous-command block-rate target.

\section{Verification targets}

\begin{center}
\begin{tabular}{@{}p{0.40\textwidth}p{0.28\textwidth}p{0.22\textwidth}@{}}
  \toprule
  \textbf{Target} & \textbf{Success criterion} & \textbf{Status}\\
  \midrule
  SO1 Fix rate & \(\geq\)80\% of failures fixed & Pending live engine\\
  SO2 Speed and cost of a fix & A full fix within 10 s, cost below 1 US cent & Pending live engine\\
  SO3 Labelling score & \(\geq\)90\% macro-F1 on the data set & Measured: 100\% clean / 99.1\% at 30\% noise\\
  SO4 Dangerous-command block rate & 100\% of dangerous commands blocked & Pending live engine\\
  \bottomrule
\end{tabular}
\end{center}

The pending targets require the full engine running on real hardware with
Docker and Ollama, which was not available when the component report was
written; their tables are intentionally left blank in the report.

\section{Public API}

The AEERE exposes a FastAPI \texttt{/execute} endpoint that accepts a
\texttt{PlannerOutput} (command and context) and returns an
\texttt{ExecutionResult} with a status of \texttt{success}, \texttt{recovered},
or \texttt{failed}, a fix log, and a failure report, together with a
\texttt{/health} check reporting whether the service and Docker are up. Sample
test cases accompany the component report.